\chapter{Perceptual Control Theory}
\label{ch:essay-pct}

The four theorems established in Chapter~\ref{ch:essay-math} describe structural properties of the admissibility relation without specifying any particular mechanism for how behavioral output $B_t$ is generated from control states or intentional structures. Perceptual Control Theory, as developed by Powers and others, offers one candidate mechanism for that generation process, though it is important to state at the outset that PCT is compatible with the admissibility framework rather than derived from it. The theorems concerning representation, coherence, repair, and non-identifiability remain valid independent of whether any particular cognitive or behavioral system operates according to control principles, and the empirical hypotheses developed in this chapter require independent validation rather than following as mathematical consequences of the formal apparatus.

The central claim of Perceptual Control Theory, as articulated by Powers, is that behavior should be understood not as a response emitted in reaction to external stimuli but as the continuously variable output of a negative-feedback process attempting to maintain selected perceptions near internally specified reference states. The fundamental control loop can be written schematically as a sequence in which a reference value $r$ is compared with a currently perceived value $p$ to produce an error signal $e = r - p$, which drives behavioral output $o$, which acts on the environment $E$, which in turn affects subsequent perception, closing the loop. What an organism controls, in this framework, is not its behavioral output directly but rather the perceptual consequences of that output. The behavior itself is therefore free to vary across a potentially wide range whenever different motor actions succeed equally well in maintaining the controlled perception near its reference value.

If linguistic behavior is generated according to control principles rather than as a fixed stimulus-response mapping, several empirical predictions follow, though each requires stating as a hypothesis rather than as an established result. First, if speakers control perceptions such as listener comprehension, social alignment, or production effort rather than controlling linguistic forms directly, then the same communicative intention should be realizable through multiple distinct surface forms whenever those forms maintain the relevant controlled variables equally well. This predicts that synonymy, paraphrase, and stylistic variation should be widespread in natural language, which is consistent with observation but does not by itself establish that control is the generating mechanism. Second, if the reference values for controlled perceptions can be adjusted dynamically in response to environmental or social conditions, then linguistic behavior should adapt to context in ways that a fixed grammar or lexicon would not predict. For instance, if ambient noise in the environment reduces the probability that phonetic distinctions are perceived correctly, a control-based account predicts that speakers should adjust their articulation to maintain comprehensibility despite the degraded channel, which could manifest as increased articulatory precision, greater reliance on redundant cues, or shifts toward more acoustically distinct phonemes. Third, if different controlled perceptions can demand conflicting adjustments, then linguistic behavior should exhibit trade-offs rather than optimizing any single dimension. A speaker might simultaneously experience pressure to increase explicitness for comprehension, reduce articulatory effort for ease of production, and conform to community norms for social alignment, with no single adjustment satisfying all three demands optimally. The resolution of such conflicts could occur through fixed hierarchical priority, dynamic adjustment of reference levels, or independent controllers operating at different linguistic levels, and each resolution mechanism predicts a different pattern of observable linguistic behavior.

The connection between control-based generation of behavior and the formal apparatus of admissibility operates through Theorem~2.1, Admissibility Preservation. Whatever mechanism generates the behavioral trajectory $B_t$, whether control-based or otherwise, the resulting trajectory through state space remains subject to the coherence condition established by that theorem. A trajectory degenerates, and thereby fails to be coherent, when the reachable set collapses to the empty set or to a singleton already determined by prior history, regardless of how the transitions were selected. A control process that drives the system into a state from which no admissible continuation exists has produced an incoherent trajectory in the technical sense defined by the theorem, even if the control loop itself remains operationally stable. The admissibility framework therefore constrains which trajectories count as coherent independent of the generative mechanism, while control theory offers one account of how a system might navigate through the space of admissible trajectories without prespecifying which particular path will be taken.

The most direct connection between PCT and the projection problem addressed in this essay concerns the equivalence class of behaviors that maintain a given controlled perception. If multiple distinct surface forms can maintain listener comprehension at the same level, those forms occupy an equivalence class with respect to the comprehension controller, even if they differ substantially in other dimensions such as phonetic detail, syntactic structure, or lexical choice. This is precisely the structure addressed by Representation Invariance: if the admissibility structure induced by the comprehension constraint is preserved across the different surface forms, then those forms count as semantically equivalent in the sense established by Theorem~1.1, despite their substrate differences. The hypothesis that natural linguistic systems operate according to control principles therefore predicts that representational equivalence classes should align with the equivalence classes defined by controlled perceptual variables, which is empirically testable by measuring whether forms judged semantically equivalent by speakers maintain the same controlled variables and whether forms that differ in controlled variables are judged semantically distinct.

The ecological examples mentioned earlier illustrate how control-based adaptation might interact with environmental constraints to produce linguistic change, though again each example requires marking as a hypothesis rather than an established mechanism. If ambient noise degrades the channel through which speech is transmitted, and if speakers control for listener comprehension rather than for preservation of particular phonetic forms, then the prediction is that sound systems should adapt toward greater robustness under the specific noise conditions encountered. In a maritime environment with wind and wave noise concentrated at low frequencies, this might favor consonant systems with more high-frequency distinctions. In an urban environment with broad-spectrum mechanical noise, it might favor vowel systems with greater formant separation or increased reliance on tonal contrasts less susceptible to spectral masking. The same environmental condition encountered by a population controlling different perceptual variables, such as articulatory ease rather than comprehensibility, would predict a different adaptive trajectory. Critically, these predictions concern the direction and character of adaptation rather than merely asserting that adaptation occurs, and they are in principle testable by comparing sound-system evolution across populations in documented acoustic environments, though such data are not yet available in the repository's implemented experiments.

The relationship between control theory and the Non-Identifiability theorem merits explicit statement. If two populations arrive at similar observable linguistic states through different control histories, one driven primarily by environmental noise favoring robust contrasts and another driven primarily by social differentiation producing superficially similar forms, then Theorem~4.2 applies: an observer examining only the contemporary languages without access to the complete control history faces a potential collision in which $\obsmap(H_1) = \obsmap(H_2)$ despite $H_1 \neq H_2$. The diagnostic corollary, Corollary~4.1, provides a criterion for assessing whether such a case represents a structural limit or a correctable observational error by checking whether the two histories share the same reachable-set structure. If the reachable sets differ, additional evidence could in principle distinguish the histories. If they are equal, no observation map factoring through reachable-set structure can resolve the ambiguity. The hypothesis that linguistic systems are generated by control processes thereby introduces a specific class of potential collisions in which different control conditions produce the same linguistic outcomes, and the theorem establishes which such collisions are fundamentally irresolvable.

These connections between control theory and the admissibility framework remain hypotheses requiring empirical validation. The theorems proven in \emph{Admissible Reconstruction} establish what follows from the admissibility relation itself, independent of any commitment to control as the generative mechanism. Control theory offers a particular account of how behavioral trajectories might be generated in a way that interacts naturally with admissibility constraints, equivalence classes, and observational limits, but that interaction is a conjecture about empirical systems rather than a mathematical consequence of the formal apparatus. The value of stating the connection explicitly is that it generates testable predictions concerning equivalence classes, environmental adaptation, conflict resolution, and historical non-identifiability that can be investigated through targeted experiments rather than remaining at the level of informal speculation.
